%----------------------------------------------------------------------------------------
% UITSPRAKEN EN WAARHEIDSTABELLEN. Alle opgaven zijn eigen werk.
%
% Opbouw: Probleem, Start, leerstof, en op een nieuwe pagina de Opdrachten en Test jezelf (geen Koppel in de 2de graad).
% Labels van dit hoofdstuk beginnen met h1-.
% Verschillen per pakket: waarheidstabellen met drie uitspraken (2-5); in de oefeningen een eigen
% variant voor pakket 1 waar de andere pakketten drie uitspraken gebruiken.
%----------------------------------------------------------------------------------------

\mychapter{1}{Uitspraken en waarheidstabellen}

\begin{probleem}
Op een eiland wonen twee soorten mensen: \textit{eerlijken}, die altijd de waarheid spreken, en \textit{leugenaars}, die altijd liegen. Aan hun uiterlijk zie je het verschil niet. Je ontmoet twee eilandbewoners, Amira en Bas. Amira zegt:

\smallskip
{\leftskip=6mm\textit{Bas en ik zijn allebei leugenaars.}\par}

\smallskip
Wie van de twee is eerlijk, en wie liegt? Kun je dat zeker weten, zonder te gokken?
\end{probleem}

\lkr{Leerplan \pakketnaam: LPD \alleenvoor{1}{26}\alleenvoor{2}{30}\alleenvoor{3}{32}\alleenvoor{4,5}{36} (dit hoofdstuk: uitspraken, de voegwoorden niet, en, of, en waarheidstabellen). De wenken bij dit doel staan erin: de exclusieve \textit{of} van de gewone taal, en schakelingen bij \textit{en} en \textit{of}.\alleenvoor{5}{ Het servicedocument voor 5 uur noemt ook: beoordelen of een zin een uitspraak is (sectie 1).} \alleenvoor{2,3,4,5}{Waarheidstabellen met drie uitspraken staan niet in pakket 1. }\alleenvoor{1}{Waarheidstabellen met drie uitspraken staan in de andere pakketten, niet in dit pakket. }Dat is een keuze van dit handboek: het leerplan is voor alle pakketten hetzelfde. Het ontkennen van \textit{alle} en \textit{er is} (sectie 4) is ook een keuze; het bereidt het tegenvoorbeeld voor (LPD 4).}

Wiskundigen spreken een precieze taal. In het dagelijks leven betekent \textit{of} soms \textit{het ene of het andere, maar niet allebei}, en soms ook \textit{allebei mag}. Je merkt het meestal niet, maar in een regel, een wet of een wiskundige eigenschap kan dat verschil groot zijn. In dit hoofdstuk leer je wat een uitspraak is, hoe je uitspraken combineert met \textit{niet}, \textit{en} en \textit{of}, en hoe je met een waarheidstabel systematisch nagaat wanneer een samengestelde uitspraak waar is. Daarmee los je ook het raadsel van het eiland op.

\begin{oefeningen}[Start]
\oef{h1-start-1}

\oef{h1-verken}
\end{oefeningen}

\lkr{Laat de Verken-oefening maken vóór sectie 2. Er komt meestal discussie over vraag b: in de taal van elke dag sluit \textit{of} soms \textit{allebei} uit, in de wiskunde nooit.}

%========================================================================================
\section{Uitspraken}

\subsection{Wat is een uitspraak?}

Kijk naar deze zinnen: \textit{7 is een priemgetal}, \textit{elke rechthoek is een vierkant}, \textit{Brussel ligt in België}. Van elk ervan kun je zeggen of ze klopt: de eerste en de derde wel, de tweede niet. Maar bij \textit{Hoe laat is het?}, \textit{Doe de deur dicht!} of \textit{Pizza is lekker} kun je dat niet zeggen. Een vraag of een bevel klopt niet en klopt ook niet fout, en over smaak verschillen mensen. In de logica werken we alleen met zinnen waarvan je kunt zeggen of ze kloppen.

\begin{theorie}
Een \textbf{\textcolor{accent}{uitspraak}} is een zin die ofwel waar ofwel onwaar is.

\vspace{1mm}
Of een uitspraak waar of onwaar is, heet haar \textbf{\textcolor{accent}{waarheidswaarde}}. We korten af: $\formule{\text{W}}$ voor waar en $\formule{\text{O}}$ voor onwaar.
\end{theorie}

Een uitspraak hoeft dus niet waar te zijn. \textit{12 is deelbaar door 5} is een uitspraak, met waarheidswaarde O. Een uitspraak is ook nooit allebei: ze is waar of ze is onwaar. Om kort te schrijven, geven we een uitspraak vaak een kleine letter, zoals $p$:~\textit{7 is een priemgetal}.

\voorbeeld{}{Welke zinnen zijn uitspraken? Geef bij elke uitspraak de waarheidswaarde.\par
\textit{(a) 15 is een veelvoud van 3.} \quad \textit{(b) Is 15 een veelvoud van 4?} \quad \textit{(c) Elke ruit is een vierkant.}\par
\textit{(d) Wiskunde is het leukste vak.} \quad \textit{(e) Teken een gelijkzijdige driehoek.}\label{vb:h1-uitspraak}}{%
\uitleg{Bij elke zin vragen we ons af: kun je zeggen of ze waar is of niet, los van wie ze leest?}
\begin{stappen}
\text{(a) uitspraak, W} & $15 = 3 \cdot 5$. \\
\text{(b) geen uitspraak} & Een vraag is waar noch onwaar. \\
\text{(c) uitspraak, O} & Een ruit met een hoek van $60^\circ$ is geen vierkant. \\
\text{(d) geen uitspraak} & Een mening: daarover verschillen mensen. \\
\text{(e) geen uitspraak} & Een opdracht is waar noch onwaar.
\end{stappen}}

\begin{oefeningen}
\oef{h1-uitspraak}
\end{oefeningen}

\subsection{Open zinnen}

En de zin $x + 3 = 10$? Of die klopt, hangt af van $x$. Vul je $7$ in, dan krijg je $7 + 3 = 10$: waar. Vul je $1$ in, dan krijg je $1 + 3 = 10$: onwaar. Zolang je niet weet wat $x$ is, kun je niet zeggen of de zin waar is.

\begin{theorie}
Een \textbf{\textcolor{accent}{open zin}} is een zin met een letter, die een uitspraak wordt als je voor die letter een waarde invult.
\end{theorie}

Ook \textit{hij is 15 jaar} is een open zin: de waarheid hangt af van wie \textit{hij} is. Een vergelijking uit de algebra is een open zin, en de oplossingen zijn net de getallen die er een ware uitspraak van maken.

\voorbeeld{}{Voor welke natuurlijke getallen $n$ van 1 tot en met 10 is de open zin \textit{$n$ is een priemgetal kleiner dan 6} waar?\label{vb:h1-openzin}}{%
\uitleg{We vullen de getallen een voor een in, en kijken telkens of we een ware uitspraak krijgen.}
\begin{stappen}
n = 2,\ 3,\ 5 & Priemgetallen en kleiner dan 6: waar. \\
n = 7 & Een priemgetal, maar niet kleiner dan 6: onwaar. \\
n = 1,\ 4,\ 6,\ 8,\ 9,\ 10 & Geen priemgetal: onwaar.
\end{stappen}
De open zin is dus waar voor $n = 2$, $3$ en $5$.}

\begin{oefeningen}
\oef{h1-openzin}
\end{oefeningen}

%========================================================================================
\section{Niet, en, of}

Met de woordjes \textit{niet}, \textit{en} en \textit{of} maak je van uitspraken nieuwe uitspraken: \textit{het regent niet}, \textit{het regent en het waait}, \textit{ik neem de fiets of de bus}. Zulke woordjes heten \textit{voegwoorden} of \textit{connectieven}. Wanneer is zo'n nieuwe uitspraak waar? Dat hangt alleen af van de delen waaruit ze bestaat.

\subsection{Niet}

Hoe zeg je dat een uitspraak niet klopt? Neem $p$:~\textit{Het regent.} Dan zegt \textit{Het regent niet} net het tegendeel. Regent het, dan is $p$ waar en de nieuwe zin onwaar. Regent het niet, dan is het omgekeerd.

\begin{theorie}
\naast{De \textbf{\textcolor{accent}{negatie}} of ontkenning van een uitspraak $p$ is de uitspraak \textit{niet $p$}, met als notatie $\formule{\niet p}$.

\vspace{1mm}
$\niet p$ is waar als $p$ onwaar is, en onwaar als $p$ waar is.
}{\begin{waarheidstabel}{2}
$p$ & $\niet p$ \\ \hline
W & O \\ \hline
O & W \\ \hline
\end{waarheidstabel}}
\end{theorie}

De tabel rechts in het kader is een \textit{waarheidstabel}: ze geeft voor elke waarheidswaarde van $p$ die van $\niet p$. Je ontkent een uitspraak door er \textit{niet} in te zetten, of door ervoor te zetten: \textit{het is niet waar dat}.

\voorbeeld{}{Ontken de uitspraak $p$:~\textit{17 is kleiner dan 20.} Geef de waarheidswaarde van $p$ en van $\niet p$.\label{vb:h1-negatie}}{%
\uitleg{We zetten \textit{niet} in de zin. Daarna zoeken we een kortere manier om hetzelfde te zeggen.}
\begin{stappen}
\niet p\colon\ \text{17 is niet kleiner dan 20} & De ontkenning, letterlijk. \\
\niet p\colon\ 17 \geqslant 20 & Niet kleiner dan 20: groter dan of gelijk aan 20. \\
p\colon \text{ W}, \quad \niet p\colon \text{ O} & $17 < 20$ klopt.
\end{stappen}
Let op: de ontkenning is niet $17 > 20$. Is een getal niet kleiner dan 20, dan kan het ook gelijk zijn aan 20.}

Zo ook in het dagelijks leven: de ontkenning van \textit{het is warm} is \textit{het is niet warm}, en niet \textit{het is koud}. Het kan ook lauw zijn.

\begin{oefeningen}
\oef{h1-negatie}
\end{oefeningen}

\subsection{En}

\naast{Een lamp hangt in een stroomkring met twee schakelaars na elkaar, $p$ en $q$ (in \textit{serie}). Wanneer brandt de lamp? Alleen als de stroom langs $p$ én langs $q$ kan, dus als beide schakelaars dicht zijn. Is er één open, dan is de kring onderbroken. Noem $p$ de uitspraak \textit{schakelaar $p$ is dicht}, en $q$ net zo. Dan brandt de lamp als $p$ en $q$ allebei waar zijn.}{\begin{tikzpicture}[scale=.62,font=\scriptsize]
  \draw[grijs,thick] (0,0) -- (1,0) (2,0) -- (3,0) (4,0) -- (4.7,0) (5.5,0) -- (6.2,0) -- (6.2,-1.3) -- (3.4,-1.3) (2.8,-1.3) -- (0,-1.3) -- (0,0);
  \draw[thick] (3.4,-1.0) -- (3.4,-1.6) (2.8,-1.15) -- (2.8,-1.45);
  \fill (1,0) circle (1.6pt) (2,0) circle (1.6pt) (3,0) circle (1.6pt) (4,0) circle (1.6pt);
  \draw[accent,very thick] (1,0) -- (1.9,.45) (3,0) -- (4,0);
  \node[above] at (1.5,.45) {$p$};
  \node[above] at (3.5,.1) {$q$};
  \draw[thick] (5.1,0) circle (.4);
  \draw (4.82,-.28) -- (5.38,.28) (4.82,.28) -- (5.38,-.28);
  \node[below] at (5.1,-.4) {lamp};
\end{tikzpicture}}

In de figuur is $q$ dicht en $p$ open: de lamp brandt niet. In gewone taal is \textit{en} net zo: de zin \textit{ik heb een fiets en ik heb een step} is alleen waar als beide delen waar zijn.

\begin{theorie}
\naast{De \textbf{\textcolor{accent}{conjunctie}} van twee uitspraken $p$ en $q$ is de uitspraak \textit{$p$ en $q$}, met als notatie $\formule{p \en q}$.

\vspace{1mm}
$p \en q$ is waar als $p$ en $q$ allebei waar zijn, en anders onwaar.
}{\begin{waarheidstabel}{3}
$p$ & $q$ & $p \en q$ \\ \hline
W & W & W \\ \hline
W & O & O \\ \hline
O & W & O \\ \hline
O & O & O \\ \hline
\end{waarheidstabel}}
\end{theorie}

Er zijn vier rijen, omdat er vier combinaties zijn: $p$ en $q$ kunnen elk waar of onwaar zijn. We zetten ze altijd in dezelfde volgorde WW, WO, OW, OO.

\voorbeeld{}{Is de uitspraak \textit{12 is even en 12 is deelbaar door 5} waar?\label{vb:h1-en}}{%
\uitleg{Het is een conjunctie. We bepalen eerst de waarheidswaarde van de twee delen, en lezen dan de rij in de tabel af.}
\begin{stappen}
p\colon \text{ 12 is even} & W, want $12 = 2 \cdot 6$. \\
q\colon \text{ 12 is deelbaar door 5} & O, want $12 : 5 = 2,4$. \\
p \en q\colon \text{ O} & De rij WO: één deel is onwaar.
\end{stappen}}

In gewone taal gebruik je voor \textit{en} ook andere woorden. \textit{Het is koud maar zonnig} zegt hetzelfde als \textit{het is koud en het is zonnig}: het woordje \textit{maar} voegt alleen een tegenstelling in de toon toe.

\begin{oefeningen}
\oef{h1-en}
\end{oefeningen}

\subsection{Of}

In de Verken-oefening mag Sara mee op kamp: ze is 12 én ze is lid. Wie de regel opstelde, bedoelde zeker niet dat wie aan beide voorwaarden voldoet, thuis moet blijven. \textit{Of} sluit \textit{allebei} daar dus niet uit.

\naast{Ook in een stroomkring zie je dat. Hier staan de schakelaars $p$ en $q$ naast elkaar (\textit{parallel}). De stroom kan langs $p$ of langs $q$ naar de lamp. De lamp brandt dus als minstens één schakelaar dicht is: alleen $p$, alleen $q$, of allebei. Alleen als beide open zijn, blijft ze uit.}{\begin{tikzpicture}[scale=.62,font=\scriptsize]
  \draw[grijs,thick] (0,0) -- (1,0) (1,.6) -- (1.5,.6) (2.5,.6) -- (3,.6) (1,-.6) -- (1.5,-.6) (2.5,-.6) -- (3,-.6) (1,.6) -- (1,-.6) (3,.6) -- (3,-.6) (3,0) -- (4.2,0) (5,0) -- (5.6,0) -- (5.6,-1.5) -- (3.1,-1.5) (2.5,-1.5) -- (0,-1.5) -- (0,0);
  \draw[thick] (3.1,-1.2) -- (3.1,-1.8) (2.5,-1.35) -- (2.5,-1.65);
  \fill (1.5,.6) circle (1.6pt) (2.5,.6) circle (1.6pt) (1.5,-.6) circle (1.6pt) (2.5,-.6) circle (1.6pt);
  \draw[accent,very thick] (1.5,.6) -- (2.4,1.05) (1.5,-.6) -- (2.5,-.6);
  \node[above] at (2,1.05) {$p$};
  \node[above] at (2,-.5) {$q$};
  \draw[thick] (4.6,0) circle (.4);
  \draw (4.32,-.28) -- (4.88,.28) (4.32,.28) -- (4.88,-.28);
\end{tikzpicture}}

\begin{theorie}
\naast{De \textbf{\textcolor{accent}{disjunctie}} van twee uitspraken $p$ en $q$ is de uitspraak \textit{$p$ of $q$}, met als notatie $\formule{p \of q}$.

\vspace{1mm}
$p \of q$ is waar als minstens één van beide waar is (ook als ze allebei waar zijn), en onwaar als ze allebei onwaar zijn.
}{\begin{waarheidstabel}{3}
$p$ & $q$ & $p \of q$ \\ \hline
W & W & W \\ \hline
W & O & W \\ \hline
O & W & W \\ \hline
O & O & O \\ \hline
\end{waarheidstabel}}
\end{theorie}

In het dagelijks leven bedoelt men met \textit{of} soms \textit{precies één van de twee}. Bij \textit{je krijgt soep of dessert} bij een menu mag je meestal niet allebei kiezen. Dat heet de \textit{exclusieve of}. In de wiskunde betekent \textit{of} altijd: het ene, het andere, of allebei. Het teken $\of$ komt waarschijnlijk van het Latijnse \textit{vel}, dat precies die betekenis had. Het teken $\en$ is $\of$ omgekeerd.

\voorbeeld{}{Is de uitspraak \textit{een vierkant is een rechthoek of een vierkant is een ruit} waar?\label{vb:h1-of}}{%
\uitleg{Een disjunctie: we bepalen de waarheidswaarde van beide delen.}
\begin{stappen}
p\colon \text{ een vierkant is een rechthoek} & W: een vierkant heeft vier rechte hoeken. \\
q\colon \text{ een vierkant is een ruit} & W: een vierkant heeft vier gelijke zijden. \\
p \of q\colon \text{ W} & De rij WW: \textit{of} sluit \textit{allebei} niet uit.
\end{stappen}}

\begin{oefeningen}
\oef{h1-of}
\end{oefeningen}

%========================================================================================
\section{Waarheidstabellen}

\subsection{Een waarheidstabel opstellen}

Met \textit{niet}, \textit{en} en \textit{of} bouw je langere uitspraken, zoals $p \en \niet q$: \textit{ik heb huiswerk en ik ga niet sporten}. Wanneer is zo'n uitspraak waar? Je kunt dat voor elk geval apart nagaan, maar dan vergeet je makkelijk een geval. Hoeveel gevallen zijn er eigenlijk? Voor $p$ zijn er twee mogelijkheden, W of O, en bij elk daarvan twee voor $q$. Dat geeft $2 \cdot 2 = 4$ gevallen, de vier rijen die je al kent. Zet je ze allemaal onder elkaar, dan vergeet je er geen.

\begin{theorie}
Een \textbf{\textcolor{accent}{waarheidstabel}} geeft de waarheidswaarde van een samengestelde uitspraak voor elke combinatie van de waarheidswaarden van haar delen. Met twee uitspraken $p$ en $q$ heeft de tabel vier rijen.
\end{theorie}

\voorbeeld{}{Maak de waarheidstabel van $p \en \niet q$.\label{vb:h1-tabel}}{%
\naast{\uitleg{We beginnen met de vier rijen voor $p$ en $q$. Om $p \en \niet q$ te kennen, hebben we eerst $\niet q$ nodig: dat krijgt een eigen kolom.}
De kolom $\niet q$ is het omgekeerde van de kolom $q$. Daarna gebruiken we de regel voor \textit{en} op de kolommen $p$ en $\niet q$: de uitkomst is alleen W als beide W zijn.

De uitspraak $p \en \niet q$ is dus alleen waar als $p$ waar is en $q$ onwaar: de rij WO.
}{\begin{waarheidstabel}{4}
$p$ & $q$ & $\niet q$ & $p \en \niet q$ \\ \hline
W & W & O & \res{O} \\ \hline
W & O & W & \res{W} \\ \hline
O & W & O & \res{O} \\ \hline
O & O & W & \res{O} \\ \hline
\end{waarheidstabel}}}

\begin{werkwijze}
Schrijf in de eerste twee kolommen de vier combinaties WW, WO, OW, OO. Maak voor elk deel van de uitspraak een kolom, van het kleinste deel naar het grootste. Bereken elke kolom met de tabel van \textit{niet}, \textit{en} of \textit{of}. De laatste kolom is de uitspraak zelf.
\end{werkwijze}

\begin{oefeningen}
\oef{h1-tabel}
\end{oefeningen}

\subsection{Haakjes}

Wat betekent $\niet p \en q$? Ontken je alleen $p$, of de hele conjunctie? We spreken af dat $\niet$ alleen hoort bij wat er net achter staat. Dus $\niet p \en q$ betekent $(\niet p) \en q$. Wil je de hele conjunctie ontkennen, dan zet je haakjes: $\niet(p \en q)$. Ook als $\en$ en $\of$ samen voorkomen, zet je altijd haakjes: $p \en q \of r$ kan $(p \en q) \of r$ of $p \en (q \of r)$ betekenen.

\voorbeeld{}{Neem $p$:~\textit{Ik neem de trein} en $q$:~\textit{Ik neem de bus}. Vergelijk $\niet p \en q$ en $\niet(p \en q)$ met een waarheidstabel. Wat betekenen ze in woorden?\label{vb:h1-haakjes}}{%
\uitleg{Bij $\niet p \en q$ hebben we eerst $\niet p$ nodig, bij $\niet(p \en q)$ eerst $p \en q$. We zetten alles in één tabel.}
\figuren{\begin{waarheidstabel}{6}
$p$ & $q$ & $\niet p$ & $\niet p \en q$ & $p \en q$ & $\niet(p \en q)$ \\ \hline
W & W & O & \res{O} & W & \res{O} \\ \hline
W & O & O & \res{O} & O & \res{W} \\ \hline
O & W & W & \res{W} & O & \res{W} \\ \hline
O & O & W & \res{O} & O & \res{W} \\ \hline
\end{waarheidstabel}}
De kolommen verschillen in twee rijen. $\niet p \en q$ betekent \textit{ik neem de trein niet en ik neem de bus}: alleen waar in de rij OW. $\niet(p \en q)$ betekent \textit{het is niet waar dat ik de trein en de bus neem}: alleen onwaar als ik ze allebei neem.}

\begin{oefeningen}
\oef{h1-haakjes}
\end{oefeningen}

\subsection{Van woorden naar symbolen en terug}

Om een uitspraak in woorden met een waarheidstabel te onderzoeken, schrijf je ze eerst in symbolen. Kies voor elke eenvoudige uitspraak een letter, en zoek dan de voegwoorden. Soms zitten die verstopt: \textit{maar} is \textit{en}, \textit{geen} en \textit{niet waar dat} zijn een negatie, en \textit{noch \ldots{} noch} betekent \textit{niet het ene en niet het andere}.

\voorbeeld{}{Neem $p$:~\textit{Ik drink koffie} en $q$:~\textit{Ik drink thee}. Schrijf in symbolen: \textit{Ik drink noch koffie, noch thee.}\label{vb:h1-vertaal}}{%
\uitleg{\textit{Noch \ldots{} noch} is een verstopte \textit{en} met twee negaties.}
\begin{stappen}
\text{ik drink geen koffie en ik drink geen thee} & Dezelfde zin, met de voegwoorden zichtbaar. \\
\niet p \en \niet q & Twee negaties, verbonden met \textit{en}.
\end{stappen}}

\voorbeeld{}{Neem dezelfde $p$ en $q$. Schrijf in woorden: $\niet(p \of q)$.\label{vb:h1-woorden}}{%
\uitleg{De negatie staat voor de haakjes: ze ontkent de hele disjunctie. Daarvoor gebruik je \textit{het is niet waar dat}.}
Eerst wat tussen de haakjes staat: $p \of q$ is \textit{ik drink koffie of thee}. Daarna de negatie ervoor:

\smallskip
\qquad $\niet(p \of q)$:~\textit{het is niet waar dat ik koffie of thee drink.}

\smallskip
Dat lijkt sterk op de zin uit het vorige voorbeeld. In sectie 4 gaan we na of ze echt hetzelfde zeggen.}

\begin{oefeningen}
\oef{h1-vertaal}
\end{oefeningen}

\begin{alleenvoor}{2,3,4,5}
\subsection{Drie uitspraken}

Het alarm van een winkel gaat af als de deur of het raam opengaat, terwijl het alarm niet uitgeschakeld is. Neem $p$:~\textit{de deur gaat open}, $q$:~\textit{het raam gaat open} en $r$:~\textit{het alarm is uitgeschakeld}. Het alarm gaat af als $(p \of q) \en \niet r$ waar is. Nu zijn er drie uitspraken, en dus $2 \cdot 2 \cdot 2 = 8$ combinaties. Om er geen te vergeten, wisselt de kolom $p$ per vier, de kolom $q$ per twee en de kolom $r$ telkens.

\voorbeeld{}{Maak de waarheidstabel van $(p \of q) \en \niet r$. Wanneer gaat het alarm af?\label{vb:h1-drie}}{%
\naast{\uitleg{Acht rijen. We berekenen eerst $p \of q$ en $\niet r$, en nemen dan de conjunctie van die twee kolommen.}
De laatste kolom is W in de rijen WWO, WOO en OWO. Het alarm gaat af als minstens één van deur en raam opengaat en het alarm niet uitgeschakeld is.

Gaat er niets open (de laatste twee rijen), dan blijft het alarm stil, ook als het aan staat.
}{\begin{waarheidstabel}{6}
$p$ & $q$ & $r$ & $p \of q$ & $\niet r$ & $(p \of q) \en \niet r$ \\ \hline
W & W & W & W & O & \res{O} \\ \hline
W & W & O & W & W & \res{W} \\ \hline
W & O & W & W & O & \res{O} \\ \hline
W & O & O & W & W & \res{W} \\ \hline
O & W & W & W & O & \res{O} \\ \hline
O & W & O & W & W & \res{W} \\ \hline
O & O & W & O & O & \res{O} \\ \hline
O & O & O & O & W & \res{O} \\ \hline
\end{waarheidstabel}}}

\begin{oefeningen}
\oef{h1-drie}
\end{oefeningen}
\end{alleenvoor}

\subsection{Terug naar het eiland}

Nu kunnen we het raadsel van het begin systematisch oplossen. Het idee: een eerlijke bewoner zegt altijd iets waars, een leugenaar altijd iets onwaars. In een rij van de waarheidstabel die kan kloppen, is wat Amira zegt dus waar als Amira eerlijk is, en onwaar als ze liegt.

\voorbeeld{}{Terug naar het probleem. Amira zegt: \textit{Bas en ik zijn allebei leugenaars.} Wie is eerlijk, en wie liegt?\label{vb:h1-terug}}{%
\uitleg{Neem $p$:~\textit{Amira is eerlijk} en $q$:~\textit{Bas is eerlijk}. Amira zegt dat geen van beiden eerlijk is: $\niet p \en \niet q$. We maken de waarheidstabel, en vergelijken in elke rij de waarheidswaarde van wat Amira zegt met die van $p$.}
\figuren{\begin{waarheidstabelbreed}{1.9cm}{6}
$p$ & $q$ & $\niet p$ & $\niet q$ & $\niet p \en \niet q$ & kan dit? \\ \hline
W & W & O & O & O & nee \\ \hline
W & O & O & W & O & nee \\ \hline
O & W & W & O & O & \res{ja} \\ \hline
O & O & W & W & W & nee \\ \hline
\end{waarheidstabelbreed}}
In de eerste twee rijen is Amira eerlijk, maar zegt ze iets onwaars: dat kan niet. In de laatste rij liegt Amira, maar zegt ze iets waars: dat kan ook niet. Alleen de derde rij klopt: Amira liegt, en wat ze zegt is inderdaad onwaar. \textbf{Amira is een leugenaar en Bas is eerlijk.}

Zonder tabel kan het ook, maar dan moet je goed opletten: was Amira eerlijk, dan zou ze zelf een leugenaar zijn. Dus liegt ze, en dan zijn ze niet allebei leugenaars: Bas is eerlijk.}

%========================================================================================
\section{Ontkennen}

\subsection{Gelijkwaardige uitspraken}

In voorbeeld \ref{vb:h1-vertaal} en \ref{vb:h1-woorden} vonden we twee zinnen die erg op elkaar lijken: \textit{ik drink noch koffie, noch thee} en \textit{het is niet waar dat ik koffie of thee drink}. Zeggen ze hetzelfde? In gewone taal is dat moeilijk te beslissen. Met een waarheidstabel kan het wel: zeggen ze hetzelfde, dan zijn ze in precies dezelfde gevallen waar.

\begin{theorie}
Twee uitspraken zijn \textbf{\textcolor{accent}{gelijkwaardig}} als ze in elke rij van de waarheidstabel dezelfde waarheidswaarde hebben.
\end{theorie}

Gelijkwaardige uitspraken mag je altijd door elkaar vervangen. Een eenvoudig voorbeeld is de dubbele ontkenning: $\niet(\niet p)$ is gelijkwaardig met $p$. \textit{Het is niet waar dat ik niet kom} betekent gewoon \textit{ik kom}.

\voorbeeld{}{Zijn $\niet(p \of q)$ en $\niet p \en \niet q$ gelijkwaardig?\label{vb:h1-gelijkw}}{%
\uitleg{We maken één tabel met een kolom voor elk deel, en vergelijken de twee kolommen die we zoeken rij per rij.}
\figuren{\begin{waarheidstabel}{7}
$p$ & $q$ & $p \of q$ & $\niet(p \of q)$ & $\niet p$ & $\niet q$ & $\niet p \en \niet q$ \\ \hline
W & W & W & \res{O} & O & O & \res{O} \\ \hline
W & O & W & \res{O} & O & W & \res{O} \\ \hline
O & W & W & \res{O} & W & O & \res{O} \\ \hline
O & O & O & \res{W} & W & W & \res{W} \\ \hline
\end{waarheidstabel}}
De twee kolommen zijn in elke rij gelijk: de uitspraken zijn gelijkwaardig. \textit{Ik drink noch koffie, noch thee} en \textit{het is niet waar dat ik koffie of thee drink} zeggen dus hetzelfde.}

\begin{werkwijze}
Maak één waarheidstabel met een kolom voor elke uitspraak. Zijn de twee kolommen in elke rij gelijk, dan zijn de uitspraken gelijkwaardig. Verschillen ze in één rij, dan zijn ze het niet.
\end{werkwijze}

\begin{oefeningen}
\oef{h1-gelijkw}
\end{oefeningen}

\subsection{De wetten van De Morgan}

Het vorige voorbeeld zegt hoe je een \textit{of} ontkent: het is niet waar dat $p$ of $q$, betekent dat $p$ niet waar is en $q$ ook niet. Die regel is genoemd naar de Britse wiskundige Augustus De Morgan (1806-1871).

\begin{theorie}
\begin{stelling}{Ontkenning van een disjunctie}
\qquad $\formule{\niet(p \of q)}$ \ is gelijkwaardig met \ $\formule{\niet p \en \niet q}$
\end{stelling}
\end{theorie}

En hoe ontken je een \textit{en}? Neem de uitspraak \textit{Mila is 15 en Mila woont in Gent}. Ze is onwaar als Mila geen 15 is, of als Mila niet in Gent woont, of allebei. Daar duikt een \textit{of} op.

\voorbeeld{}{Zijn $\niet(p \en q)$ en $\niet p \of \niet q$ gelijkwaardig?\label{vb:h1-morgan}}{%
\uitleg{Zoals in het vorige voorbeeld: één tabel, en de twee kolommen vergelijken.}
\figuren{\begin{waarheidstabel}{7}
$p$ & $q$ & $p \en q$ & $\niet(p \en q)$ & $\niet p$ & $\niet q$ & $\niet p \of \niet q$ \\ \hline
W & W & W & \res{O} & O & O & \res{O} \\ \hline
W & O & O & \res{W} & O & W & \res{W} \\ \hline
O & W & O & \res{W} & W & O & \res{W} \\ \hline
O & O & O & \res{W} & W & W & \res{W} \\ \hline
\end{waarheidstabel}}
De kolommen zijn gelijk: de uitspraken zijn gelijkwaardig.}

\begin{theorie}
\begin{stelling}{Ontkenning van een conjunctie}
\qquad $\formule{\niet(p \en q)}$ \ is gelijkwaardig met \ $\formule{\niet p \of \niet q}$
\end{stelling}
\end{theorie}

Kort gezegd: bij het ontkennen ontken je elk deel, en vervang je \textit{en} door \textit{of} en \textit{of} door \textit{en}. Een veelgemaakte fout is het voegwoord te laten staan. De ontkenning van \textit{Mila is 15 en woont in Gent} is niet \textit{Mila is geen 15 en woont niet in Gent}.

\voorbeeld{}{Ontken de open zin $x > 2 \ \en \ x < 5$, en schrijf de ontkenning zonder het woord \textit{niet}.\label{vb:h1-morgan2}}{%
\uitleg{Een conjunctie: met De Morgan wordt het een disjunctie van de twee ontkenningen.}
\begin{stappen}
\niet(x > 2) \ \of \ \niet(x < 5) & De ontkenning van een conjunctie. \\
x \leqslant 2 \ \of \ x \geqslant 5 & Niet groter dan 2 is kleiner dan of gelijk aan 2.
\end{stappen}
De open zin zegt dat $x$ tussen 2 en 5 ligt. De ontkenning zegt dat $x$ links van 2 of rechts van 5 ligt (of erop). Een getal kan niet tegelijk $\leqslant 2$ en $\geqslant 5$ zijn: daarom zou \textit{en} hier fout zijn.}

\begin{oefeningen}
\oef{h1-morgan}
\end{oefeningen}

\subsection{Ontkennen met alle en er is}

De leerkracht zegt: \textit{Alle leerlingen hebben hun taak gemaakt.} Lies weet dat dat niet klopt. Wat weet ze dan? Niet dat niemand zijn taak maakte. Wel dat er minstens één leerling is die zijn taak niet maakte. Eén leerling zonder taak is genoeg om de uitspraak van de leerkracht onwaar te maken.

\begin{theorie}
\begin{stelling}{Ontkenning van alle}
De ontkenning van \ $\formule{\text{alle \ldots{} zijn \ldots}}$ \ is \ $\formule{\text{er is minstens één \ldots{} dat niet \ldots{} is}}$.
\end{stelling}
\end{theorie}

\voorbeeld{}{Ontken de uitspraak \textit{Alle priemgetallen zijn oneven.} Welke van de twee is waar?\label{vb:h1-alle}}{%
\uitleg{Volgens de regel zoeken we niet \textit{geen enkel priemgetal is oneven}, maar \textit{er is een priemgetal dat niet oneven is}.}
\begin{stappen}
\text{Er is een priemgetal dat even is.} & De ontkenning: niet oneven is even. \\
2 \text{ is een even priemgetal} & De ontkenning is waar. \\
\text{De uitspraak zelf is onwaar.} & Eén geval volstaat.
\end{stappen}}

Omgekeerd ontken je \textit{er is} met \textit{geen enkel}. De uitspraak \textit{er is een leerling die zijn taak vergat} is onwaar als geen enkele leerling zijn taak vergat, dus als alle leerlingen hun taak bij hadden.

\begin{theorie}
\begin{stelling}{Ontkenning van er is}
De ontkenning van \ $\formule{\text{er is een \ldots{} dat \ldots{} is}}$ \ is \ $\formule{\text{geen enkel \ldots{} is \ldots}}$, \ of ook \ $\formule{\text{alle \ldots{} zijn niet \ldots}}$.
\end{stelling}
\end{theorie}

\voorbeeld{}{Ontken de uitspraak \textit{Er is een reëel getal waarvan het kwadraat negatief is.}\label{vb:h1-eris}}{%
\uitleg{De uitspraak begint met \textit{er is}: de ontkenning begint met \textit{geen enkel}.}
De ontkenning is \textit{Geen enkel reëel getal heeft een negatief kwadraat.} Met \textit{alle} zeg je hetzelfde: \textit{Alle reële getallen hebben een kwadraat dat niet negatief is.}

Hier is de ontkenning waar: een kwadraat is nooit negatief.}

\begin{oefeningen}
\oef{h1-alle}
\end{oefeningen}

%========================================================================================
\opdrachtenkop

\begin{oefeningen}
\oef{h1-o1}

\oef{h1-o3}

\oef{h1-o2}

\oef{h1-o6}

\oef{h1-o4}

\oef{h1-o8}

\oef{h1-o5}

\oef{h1-o7}

\oef{h1-o9}

\oef{h1-o10}

\oef{h1-o14}

\oef{h1-o11}

\oef{h1-o18}

\oef{h1-o13}

\oef{h1-o12}

\oef{h1-o15}

\oef{h1-o19}

\oef{h1-o16}

\oef{h1-o17}

\oef{h1-o21}

\oef{h1-o20}

\oef{h1-o22}
\end{oefeningen}

%----------------------------------------------------------------------------------------
\opdrachtsectie{Test jezelf}

\begin{oefeningen}
\oef{h1-t1}

\oef{h1-t2}

\oef{h1-t3}

\oef{h1-t4}

\oef{h1-t5}

\oef{h1-t6}

\oef{h1-t7}

\oef{h1-t8}
\end{oefeningen}

\herhaalschema{%
  \fv{oef:h1-t1}{vb:h1-uitspraak}{oef:h1-uitspraak}%
  \fv{oef:h1-t2}{vb:h1-negatie}{oef:h1-negatie}%
  \fv{oef:h1-t3}{vb:h1-en}{oef:h1-o4}%
  \fv{oef:h1-t4}{vb:h1-tabel}{oef:h1-tabel}%
  \fv{oef:h1-t5}{vb:h1-vertaal}{oef:h1-vertaal}%
  \fv{oef:h1-t6}{vb:h1-morgan2}{oef:h1-morgan}%
  \fv{oef:h1-t7}{vb:h1-alle}{oef:h1-alle}%
  \alleenvoor{1}{\fv{oef:h1-t8}{vb:h1-terug}{oef:h1-o14}}%
  \alleenvoor{2,3,4,5}{\fv{oef:h1-t8}{vb:h1-drie}{oef:h1-drie}}}

\ifzelfstudie
\section*{Maak je eigen samenvatting}
Sluit dit boek en schrijf in je schrift: wat een uitspraak en een open zin zijn, de waarheidstabellen van niet, en en of, hoe je een waarheidstabel opstelt, wanneer twee uitspraken gelijkwaardig zijn, de wetten van De Morgan, en hoe je \textit{alle} en \textit{er is} ontkent. Vergelijk daarna met de kaders.
\fi
